% Real-Time Sign Language Sentence Recognition — Springer LNCS
% Version: improved full draft with Cosine vs Seq2Seq comparison
%
\documentclass[runningheads]{llncs}
%
\usepackage[T1]{fontenc}
\usepackage{graphicx}
\usepackage{amsmath,amssymb}
\usepackage{booktabs}
%
% Uncomment the following two lines for blue hyperlinks (Springer eBook style):
%\usepackage{color}
%\renewcommand\UrlFont{\color{blue}\rmfamily}
%
\begin{document}
%
\title{Real-Time Sign Language Sentence Recognition:\\
Comparing Cosine Retrieval and Seq2Seq Attention Approaches}
%
%\titlerunning{Real-Time Sign Language Sentence Recognition}
%
\author{Gulshan Sariyeva\inst{1} \and
Rahima Karimova\inst{1} \and
Vasila Aliyeva\inst{1} \and
Ilaha Jamalli\inst{1}}
%
\authorrunning{G.\ Sariyeva et al.}
%
\institute{
School of IT and Engineering, ADA University, Baku, Azerbaijan\\
\email{\{gsariyeva18090, rkarimova17347, valiyeva, ijamalli\}@ada.edu.az}
}
%
\maketitle
%
\begin{abstract}
Sign language is the primary means of communication for Deaf and
Hard-of-Hearing (DHH) communities, yet automated sentence-level
recognition remains computationally expensive and data-hungry.
Low-resource sign languages such as Azerbaijani Sign Language (AzSL)
lack both large annotated corpora and lightweight recognition models.
We investigate two complementary approaches for sentence-level AzSL
recognition using a crowdsourced Telegram
dataset~\cite{mustafazada2025crowdsourcing} comprising 335~videos
across 120~sentences (119~translator, 216~user recordings).
The first is a \emph{training-free, retrieval-based pipeline} that
extracts MediaPipe hand landmarks (21~joints $\times$ 3~coordinates
$\times$ 2~hands $=$ 126 dimensions per frame), uniformly samples
64~frames from each video, constructs a compact
$64 \times 126$ feature matrix, and retrieves the closest matching
sentence via cosine similarity.
The second is a \emph{sequence-to-sequence (Seq2Seq) neural model}
that encodes video frames through a SqueezeNet~CNN followed by a
bidirectional LSTM, and decodes sentence tokens with an
attention-augmented LSTM decoder using beam search.
A real-time webcam module detects signing onset and offset through a
four-state finite state machine and returns the top-3 candidate
sentences.
The Seq2Seq pipeline achieves 35.5\% top-1 accuracy (119/335
correct), with a corpus-average BLEU-4 of 31.43\% and WER of
64.05\%.
User videos (37.0\% top-1) slightly outperform translator
videos (32.8\%), demonstrating that crowdsourced data from diverse
signers can support learned models.
Both approaches maintain sub-20\,ms matching latency on a
standard CPU.
The results highlight challenges and opportunities for
neural sign-language recognition in low-resource, crowdsourced
settings.

\keywords{Sign Language Recognition \and MediaPipe \and
Hand Landmarks \and Cosine Similarity \and LSTM \and
Seq2Seq Attention \and Real-Time Inference \and
Azerbaijani Sign Language \and Feature Matching \and
Crowdsourcing}
\end{abstract}
%
%
% ================================================================
\section{Introduction}
% ================================================================

Sign language recognition (SLR) has emerged as a crucial technology for
bridging the communication gap between Deaf and Hard-of-Hearing (DHH)
communities and hearing individuals.
While recent deep-learning systems have made significant advances on
isolated sign and word recognition, real-time \emph{sentence-level}
recognition remains particularly challenging due to variability in
signing speed, signer style, camera angle, and environmental
conditions~\cite{liu2024skeleton}.

Traditional approaches rely on handcrafted features and rule-based
systems that often struggle to generalise across signers and
contexts~\cite{fang2013bayesian}.
More recent work leverages skeleton-based representations and
spatio-temporal modelling to capture hand-motion dynamics and
inter-joint correlations, but these methods typically require extensive
training data and GPU resources.
Low-resource languages such as Azerbaijani Sign Language (AzSL) are
especially underserved: the Azerbaijani Sign Language Dataset
(AzSLD)~\cite{alishzade2024azsld} provides an important starting
point, yet lightweight recognition systems for AzSL—whether
training-free or learned—remain largely absent.

We address this gap by investigating and comparing two complementary
approaches on a crowdsourced Telegram
dataset~\cite{mustafazada2025crowdsourcing} of 335~AzSL
sentence-level videos:
\begin{enumerate}
    \item A \textbf{retrieval-based pipeline} that combines MediaPipe
          hand landmarks, wrist-centred normalisation, and cosine
          similarity matching.
          Each video is reduced to a fixed-length feature matrix,
          flattened into a global embedding, and compared against a
          pre-computed database of sentence templates.
    \item A \textbf{Seq2Seq neural pipeline} that extracts visual
          features via a SqueezeNet CNN, encodes them with a
          bidirectional LSTM, and decodes target sentences token by
          token using an attention-augmented LSTM decoder with
          beam search.
\end{enumerate}

A shared webcam-based inference module detects signing onset and offset
automatically, returning the top-3 candidate matches in real time on a
CPU.

Our contributions are as follows:
\begin{enumerate}
    \item An end-to-end, \textbf{training-free cosine retrieval}
          pipeline: video $\rightarrow$ landmarks $\rightarrow$ cosine
          match, requiring no GPU and no labelled training set.
    \item A \textbf{Seq2Seq encoder--decoder} pipeline with SqueezeNet
          visual features, bidirectional LSTM encoding, and Bahdanau
          attention decoding, adapted for sentence-level sign language
          translation.
    \item A \textbf{systematic comparison} of both approaches on
          crowdsourced Azerbaijani Sign Language sentences, with
          detailed analysis of sequence-level metrics (BLEU-4, WER,
          CER), signer-type stratification, and mode collapse
          behaviour.
    \item A \textbf{real-time webcam inference} system with automatic
          sign segmentation via a four-state finite state machine.
    \item Analysis of frame-count sensitivity, latency breakdown, and
          training dynamics over 71~epochs.
\end{enumerate}

The remainder of the paper is organised as follows.
Section~\ref{sec:related} surveys related work.
Section~\ref{sec:method} details the two proposed methods.
Section~\ref{sec:experiments} presents experiments and comparative
results, and Section~\ref{sec:conclusion} concludes with future
directions.

% ================================================================
\section{Related Work}\label{sec:related}
% ================================================================

\paragraph{Sign Language Recognition.}
Early SLR systems focused on isolated word or fingerspelling recognition
using Hidden Markov Models (HMMs) or Dynamic Time Warping (DTW) for
temporal alignment~\cite{fang2013bayesian}.
Recent approaches employ convolutional and recurrent neural
networks~\cite{liu2024skeleton}, graph convolutional
networks, and transformers to model spatio-temporal
dependencies.
While these models achieve high accuracy on benchmark datasets, they
require large labelled corpora and GPU resources, limiting their
applicability to low-resource languages.

\paragraph{Hand Pose Estimation and MediaPipe.}
Skeleton-based representations have gained popularity because they
are robust to lighting and background variations.
MediaPipe~Hands provides a lightweight, real-time pipeline for
hand-landmark estimation, outputting 21 three-dimensional joint
coordinates per hand on CPU.
Compared to OpenPose, MediaPipe offers comparable accuracy with
significantly lower computational cost, making it suitable for
edge-device deployment.

\paragraph{Sequence-to-Sequence Models for SLR.}
Encoder--decoder architectures with attention have been widely
adopted for sign language translation, where the goal is to map a
sequence of video frames to a natural-language
sentence~\cite{camgoz2018neural}.
The encoder typically consists of a CNN backbone (e.g.\ ResNet,
SqueezeNet) followed by an RNN (LSTM or GRU) to capture temporal
dynamics, while the decoder generates output tokens conditioned on
attention-weighted encoder states.
Beam search is commonly used at inference time to explore multiple
hypotheses.
However, these models require substantial training data to learn
robust representations, and performance degrades significantly in
low-resource settings.

\paragraph{Similarity-Based and Metric-Learning Approaches.}
Cosine similarity has been widely used in speaker verification and
metric learning to compare high-dimensional
embeddings~\cite{dehak2010cosine,zhang2020deep,zhu2020orthogonality}.
Retrieval-based methods, which match an input gesture to a pre-recorded
database using feature vectors, eliminate the need for training while
maintaining interpretable and scalable
inference~\cite{draganov2024hidden,weber2025timepoint}.
However, no prior work applies landmark--cosine retrieval at the
\emph{sentence level} for Azerbaijani Sign Language—nor compares it
against a trained Seq2Seq model—a gap our work addresses.

% ================================================================
\section{Proposed Methods}\label{sec:method}
% ================================================================

We present two approaches for sentence-level AzSL recognition:
(1)~a training-free cosine retrieval pipeline based on MediaPipe hand
landmarks, and (2)~a Seq2Seq encoder--decoder model with CNN visual
features and attention.
Both share a common real-time webcam inference module.

\begin{figure}[t]
\centering
\includegraphics[width=\textwidth]{diagram/systemarchaidt.drawio.png}
\caption{System architecture overview.
Offline: reference videos are processed into feature representations
and stored.
Online: webcam frames undergo the same pipeline and are matched via
cosine retrieval or decoded by the Seq2Seq model.}
\label{fig:architecture}
\end{figure}

% ================================================================
\subsection{Approach~1: Cosine Similarity Retrieval}\label{sec:cosine_method}
% ================================================================

% ----------------------------------------------------------------
\subsubsection{Frame Extraction.}\label{sec:frame}
% ----------------------------------------------------------------

To account for varying signing speeds and frame rates across devices, we
implement uniform temporal sampling using OpenCV.
Each video sequence of length~$T$ is reduced to a fixed sequence of
$N = 64$ frames, ensuring temporal alignment across the database.
The sampling index for each frame~$i$ is defined as:
\begin{equation}
\mathrm{idx}_i = \left\lfloor \frac{i \cdot T}{N} \right\rfloor,
\quad i = 0, 1, \dots, N-1\,.
\label{eq:sampling}
\end{equation}
When $T < N$, the system employs a zero-order hold (frame duplication)
to maintain the gesture's motion profile without introducing artefacts.
For $T > N$, stride-based sampling selects representative frames,
ensuring a consistent temporal dimension for the feature matrix.

% ----------------------------------------------------------------
\subsubsection{Hand Landmark Feature Extraction.}\label{sec:features}
% ----------------------------------------------------------------

Feature extraction utilises MediaPipe~Hands to generate a
126-dimensional vector per frame
($21~\text{landmarks} \times 3~\text{coordinates}\;(x,y,z) \times
2~\text{hands}$).
To ensure invariance to camera scale and signer distance—a critical
requirement for edge deployment—we apply wrist-centred normalisation.

Let $\mathbf{l}_i \in \mathbb{R}^{3}$ denote the $i$-th landmark
coordinate and $\mathbf{l}_0$ the wrist landmark.
The coordinates are first centred on the wrist and then scaled to a unit
sphere:
\begin{equation}
\mathbf{c}_i = \mathbf{l}_i - \mathbf{l}_0
\label{eq:centre}
\end{equation}
\begin{equation}
\mathbf{n}_i = \frac{\mathbf{c}_i}
                     {\max_{j}\, \lVert \mathbf{c}_j \rVert}
\label{eq:normalise}
\end{equation}
Centring on the wrist eliminates variance caused by the signer's
position within the frame, while scaling by the maximum landmark
distance ensures the gesture magnitude remains consistent regardless
of signer-to-camera distance.
When a hand is not detected, zero-padding preserves the dimensionality
of the feature vector.

\begin{figure}[t]
\centering
\includegraphics[width=0.45\linewidth]{diagram/Screenshot 2026-02-12 233701.png}
\hfill
\includegraphics[width=0.45\linewidth]{diagram/fig2_mediapipe_hand.png}
\caption{MediaPipe hand landmark model: 21 joints with $(x,y,z)$
coordinates, yielding 63~features per hand.}
\label{fig:mediapipe}
\end{figure}

% ----------------------------------------------------------------
\subsubsection{Feature Matrix and Cosine Similarity.}\label{sec:cosine}
% ----------------------------------------------------------------

The sampled sequence is structured as a feature matrix
$\mathbf{F} \in \mathbb{R}^{64 \times 126}$, which is flattened into
a global embedding $\mathbf{f} \in \mathbb{R}^{8064}$.
Inference is performed by computing the cosine similarity between the
input embedding~$\mathbf{a}$ and each database
template~$\mathbf{b}$:
\begin{equation}
\mathrm{sim}(\mathbf{a}, \mathbf{b}) =
  \frac{\mathbf{a} \cdot \mathbf{b}}
       {\lVert \mathbf{a} \rVert \; \lVert \mathbf{b} \rVert}\,.
\label{eq:cosine}
\end{equation}
The cosine metric treats each embedding as a \emph{direction} in
$\mathbb{R}^{8064}$, measuring angular proximity independently of
vector magnitude.
Because wrist-centred normalisation already constrains each frame's
coordinates to a unit sphere, the resulting embeddings naturally occupy
a bounded region of the space, mitigating the
\emph{embedding-norm effect} observed in similarity-based
optimisation~\cite{draganov2024hidden}.

For each query video, the per-sentence score is computed as the
\emph{maximum} cosine similarity among all reference videos of that
sentence.
A limitation of this flat-embedding approach is that adjacent fingers
tend to move together, introducing feature correlation that may reduce
discriminative power for visually similar sentences.

% ================================================================
\subsection{Approach~2: Seq2Seq Encoder--Decoder}\label{sec:seq2seq_method}
% ================================================================

The second approach treats sign-language sentence recognition as a
sequence-to-sequence translation task, mapping a variable-length
sequence of video frames to a variable-length token sequence
representing the target sentence.
The architecture comprises three stages: visual feature extraction,
temporal encoding, and attentive decoding.

% ----------------------------------------------------------------
\subsubsection{Visual Feature Extraction.}\label{sec:cnn_features}
% ----------------------------------------------------------------

Each video frame is first processed by a hand-detection module
(MediaPipe Hands) that crops and retains only frames containing
visible hand landmarks.
The retained frames are resized to $224 \times 224$ pixels and
passed through a pre-trained SqueezeNet~1.1
backbone~\cite{iandola2016squeezenet}, whose final feature map
($512 \times 13 \times 13 = 86{,}528$ dimensions) captures
high-level visual representations of hand shape and pose.
SqueezeNet is chosen for its compact architecture
($\sim$1.2\,M parameters), enabling CPU-friendly deployment while
retaining discriminative visual features.

% ----------------------------------------------------------------
\subsubsection{Bidirectional LSTM Encoder.}\label{sec:encoder}
% ----------------------------------------------------------------

The CNN features are fed sequentially into a bidirectional LSTM
encoder with hidden size $h = 256$:
\begin{equation}
\overrightarrow{\mathbf{h}_t} = \mathrm{LSTM}_{\rightarrow}
  (\mathbf{x}_t, \overrightarrow{\mathbf{h}_{t-1}}), \quad
\overleftarrow{\mathbf{h}_t} = \mathrm{LSTM}_{\leftarrow}
  (\mathbf{x}_t, \overleftarrow{\mathbf{h}_{t+1}})
\label{eq:bilstm}
\end{equation}
\begin{equation}
\mathbf{h}_t = [\overrightarrow{\mathbf{h}_t} \;;\;
                 \overleftarrow{\mathbf{h}_t}]
  \in \mathbb{R}^{512}
\label{eq:concat}
\end{equation}
where $\mathbf{x}_t \in \mathbb{R}^{86528}$ is the flattened CNN
output at frame~$t$, and $[\cdot\,;\,\cdot]$ denotes concatenation.
The bidirectional architecture captures both forward and backward
temporal context, producing a 512-dimensional hidden state per
frame.
The encoder outputs serve as the memory for attention-based
decoding.

% ----------------------------------------------------------------
\subsubsection{Attention Decoder.}\label{sec:decoder}
% ----------------------------------------------------------------

The decoder is an LSTM with Bahdanau-style (additive) attention over
the encoder outputs.
At each decoding step~$s$, the decoder computes an attention
distribution over the encoder states:
\begin{equation}
e_{s,t} = \mathbf{v}^\top \tanh\!\bigl(
  \mathbf{W}_q \,[\mathbf{e}_s \;;\; \mathbf{d}_{s-1}]
  + \mathbf{W}_k \,\mathbf{h}_t
\bigr)
\label{eq:energy}
\end{equation}
\begin{equation}
\alpha_{s,t} = \frac{\exp(e_{s,t}\,/\,\tau)}
                     {\sum_{t'} \exp(e_{s,t'}\,/\,\tau)}
\label{eq:attention}
\end{equation}
where $\mathbf{e}_s$ is the embedding of the previous token,
$\mathbf{d}_{s-1}$ is the decoder hidden state,
$\mathbf{W}_q$, $\mathbf{W}_k$, $\mathbf{v}$ are learnable
parameters, and $\tau$ is a temperature hyperparameter that
controls the sharpness of the attention distribution (we use
$\tau = 2.0$ to prevent attention collapse in early training
epochs).
Attention dropout ($p = 0.1$) is applied during training to
regularise the attention weights.

The context vector $\mathbf{c}_s = \sum_t \alpha_{s,t}\,\mathbf{h}_t$
is concatenated with the embedded input and passed through the decoder
LSTM:
\begin{equation}
\mathbf{d}_s = \mathrm{LSTM}(\mathrm{ReLU}(
  \mathbf{W}_{\text{attn}}\,[\mathbf{e}_s \;;\; \mathbf{c}_s]),
  \,\mathbf{d}_{s-1})
\label{eq:decoder_step}
\end{equation}
The output is projected to the vocabulary space via a linear layer
followed by log-softmax.

% ----------------------------------------------------------------
\subsubsection{Training.}\label{sec:training}
% ----------------------------------------------------------------

The model is trained with cross-entropy loss and the Adam optimiser
(initial learning rate $\eta = 0.01$).
We employ:
\begin{itemize}
    \item \textbf{Scheduled sampling}: teacher forcing ratio starts at
          50\% and decays linearly over epochs, encouraging the model to
          rely on its own predictions during training.
    \item \textbf{Learning rate scheduling}: ReduceLROnPlateau with
          factor~0.5 and patience~3.
    \item \textbf{Early stopping}: patience of 7~epochs with minimum
          improvement $\delta = 0.001$ on validation loss.
    \item \textbf{Gradient clipping}: maximum norm of~5.0.
\end{itemize}
The best model (lowest validation loss) is retained.
Training converges within 71~epochs on our dataset (best validation
loss 2.568 at epoch~64), with learning rate reduced four times via
ReduceLROnPlateau.

% ----------------------------------------------------------------
\subsubsection{Beam Search Decoding.}\label{sec:beam}
% ----------------------------------------------------------------

At inference time, beam search with width $B = 5$ is used to explore
multiple decoding hypotheses.
For each input video, the top-$K$ ($K \leq 5$) unique sentence
hypotheses are returned, ranked by cumulative log-probability.
This enables Top-$K$ accuracy evaluation analogous to the cosine
retrieval pipeline.

% ================================================================
\subsection{Real-Time Inference}\label{sec:realtime}
% ================================================================

Both approaches share a common real-time webcam inference module.
Fluid interaction is managed via a four-state finite state machine:
\begin{center}
\texttt{IDLE} $\;\rightarrow\;$ \texttt{RECORDING} $\;\rightarrow\;$
\texttt{COOLDOWN} $\;\rightarrow\;$ \texttt{MATCHING}
\end{center}

\begin{itemize}
    \item \textbf{IDLE}: The system waits for MediaPipe to detect hand
          landmarks in the webcam feed.
    \item \textbf{RECORDING}: Upon hand detection, frames are
          continuously recorded.
    \item \textbf{COOLDOWN}: When hands disappear, a 2-second temporal
          buffer begins.  If hands reappear within this window, recording
          resumes; otherwise the gesture is deemed complete.
    \item \textbf{MATCHING}: The recorded frames are trimmed (removing
          leading/trailing frames without hands), uniformly sampled to
          $N = 64$, processed through the selected pipeline (cosine
          retrieval or Seq2Seq decoding), and the top-3 results are
          displayed to the user.
\end{itemize}

The 2-second cooldown prevents premature matching of partial sentences
and balances technical accuracy with user experience.
After matching, the system returns to \texttt{IDLE}, ready for the next
sign.

% ================================================================
\section{Experiments and Results}\label{sec:experiments}
% ================================================================

This section presents the empirical evaluation of both pipelines.
We report dataset characteristics, system configuration, recognition
accuracy, comparative analysis, frame-count sensitivity, and
end-to-end inference latency.

% ----------------------------------------------------------------
\subsection{Dataset}\label{sec:dataset}
% ----------------------------------------------------------------

We evaluate on a curated set of Azerbaijani Sign Language (AzSL)
sentence-level videos collected via a Telegram bot crowdsourcing
platform~\cite{mustafazada2025crowdsourcing}, where community
members record sign language videos in their own environments.
The broader AzSLD corpus~\cite{alishzade2024azsld} provides valuable
context for the Azerbaijani DHH community; our work operates on
sentence-level subsets tailored for retrieval-based and generative
evaluation.

Videos are organised by sentence name (120~unique sentences) and
exhibit high variability in lighting, background, camera angle, and
signing style—reflecting realistic deployment conditions.
The dataset contains recordings from both professional translators
and community users who contributed through the Telegram bot.

Table~\ref{tab:dataset} summarises the dataset.

\begin{table}[t]
\centering
\caption{Telegram crowdsourced dataset statistics.}
\label{tab:dataset}
\begin{tabular}{@{}lr@{}}
\toprule
Metric & Value \\
\midrule
Total videos                & 335 \\
Unique sentences            & 120 \\
Translator videos           & 119 (35.5\%) \\
User (crowdsourced) videos  & 216 (64.5\%) \\
Videos / sentence (avg)     & 2.8 \\
Videos / sentence (min--max)& 1--7 \\
Sentences with 1 video      & 26 \\
Sentences with $\geq$2 videos & 94 \\
Avg sentence length         & 4.9 words \\
Unique vocabulary tokens    & 356 \\
Collection method           & Crowdsourced (Telegram bot) \\
Language                    & Azerbaijani SL \\
\bottomrule
\end{tabular}
\end{table}

The sentences cover everyday phrases, greetings, cultural expressions,
and descriptions commonly used by the Azerbaijani DHH community.
The mix of translator and user recordings introduces realistic
intra-class variance in signing tempo and hand articulation.

% ----------------------------------------------------------------
\subsection{Implementation Details}\label{sec:implementation}
% ----------------------------------------------------------------

Table~\ref{tab:system} lists the system hyperparameters for both
approaches.
Both pipelines run entirely on CPU without GPU acceleration.

\begin{table}[t]
\centering
\caption{System parameters for both approaches.}
\label{tab:system}
\begin{tabular}{@{}lll@{}}
\toprule
Parameter & Cosine Retrieval & Seq2Seq \\
\midrule
Sampled frames ($N$) & 64 & Variable \\
Features per frame & 126 (MediaPipe) & 86{,}528 (SqueezeNet) \\
Embedding dimension & 8{,}064 & 512 (biLSTM output) \\
Encoder hidden size & --- & 256 ($\times 2$ bidir) \\
Decoder hidden size & --- & 512 \\
Vocabulary size & --- & 356 tokens \\
Detection confidence & 0.5 & 0.5 \\
Normalisation & Wrist-centred & ImageNet stats \\
Similarity/decoding & Cosine similarity & Beam search ($B\!=\!5$) \\
Training required & No & Yes (71 epochs) \\
Cooldown timeout & 2.0\,s & 2.0\,s \\
\bottomrule
\end{tabular}
\end{table}

\noindent
\textbf{Software:} Python~3.11, MediaPipe~0.10.x, PyTorch, OpenCV,
NumPy, SciPy.\\
\textbf{Hardware:} Intel Core~i7 (13th~Gen), 16\,GB RAM—no GPU.\\
\textbf{Cosine protocol:} Leave-one-out cross-validation; each of the
335~videos serves as a query against the remaining entries.
Per-sentence scores are computed as the maximum cosine similarity
among all reference videos of that sentence.\\
\textbf{Seq2Seq protocol:} The model is trained on 335~videos
(80/20 train--validation split) for 71~epochs with early stopping.
Beam search ($B = 5$) returns the top-5 hypotheses per video.

% ----------------------------------------------------------------
\subsection{Comparative Accuracy}\label{sec:accuracy}
% ----------------------------------------------------------------

Table~\ref{tab:comparison} presents the central comparison between
cosine retrieval and the Seq2Seq model on the Telegram dataset.

\begin{table}[t]
\centering
\caption{Recognition accuracy: Cosine Retrieval vs.\ Seq2Seq
Encoder--Decoder on the Telegram crowdsourced dataset.}
\label{tab:comparison}
\begin{tabular}{@{}lccc@{}}
\toprule
Model & Videos & Training & Top-1 (\%) \\
\midrule
Cosine Retrieval      & 335 & None      & 30.7 \\
Seq2Seq + Attention   & 335 & 71 epochs & \textbf{35.5} \\
\bottomrule
\end{tabular}
\end{table}

The Seq2Seq model outperforms the training-free cosine retrieval
approach by 4.8 percentage points on Top-1 accuracy.
Despite the small dataset size, the learned encoder--decoder
benefits from the diversity of crowdsourced signer styles, enabling
it to extract generalisable visual-to-linguistic mappings.
Table~\ref{tab:seq2seq_metrics} provides additional sequence-level
metrics.

\begin{table}[t]
\centering
\caption{Seq2Seq sequence-level metrics on the Telegram dataset.}
\label{tab:seq2seq_metrics}
\begin{tabular}{@{}lrrr@{}}
\toprule
Metric & Overall & Translator & User \\
\midrule
Top-1 Accuracy (\%) & 35.5 & 32.8 & 37.0 \\
BLEU-4 (\%)         & 31.43 & 31.09 & 31.62 \\
WER (\%)            & 64.05 & 65.94 & 63.00 \\
CER (\%)            & 50.64 & 52.88 & 49.41 \\
\bottomrule
\end{tabular}
\end{table}

% ----------------------------------------------------------------
\subsection{Telegram Seq2Seq: Signer-Type Stratification}\label{sec:telegram_signer}
% ----------------------------------------------------------------

Table~\ref{tab:telegram_detail} reports the Seq2Seq Top-1 accuracy
on the Telegram dataset, stratified by signer type.

\begin{table}[t]
\centering
\caption{Seq2Seq accuracy stratified by signer type (Telegram dataset).}
\label{tab:telegram_detail}
\begin{tabular}{@{}lcc@{}}
\toprule
Signer Type & Videos & Top-1 (\%) \\
\midrule
Translator ($n{=}119$) & 119 & 32.8 \\
User ($n{=}216$) & 216 & 37.0 \\
\midrule
All ($n{=}335$) & 335 & \textbf{35.5} \\
\bottomrule
\end{tabular}
\end{table}

User videos achieve 4.2 percentage points higher accuracy than
translator videos.
The larger number of user videos (216 vs.\ 119) provides denser
per-sentence coverage, improving the model's ability to learn
generalisable representations.
This pattern is consistent with the cosine retrieval observations,
where user-contributed videos also yield higher similarity scores
due to richer intra-class sampling.

% ----------------------------------------------------------------
\subsection{Cosine Retrieval: Signer-Type Stratification}\label{sec:signer}
% ----------------------------------------------------------------

Table~\ref{tab:cosine_detail} reports Top-$K$ accuracy for the
cosine retrieval approach on the Telegram dataset, stratified by
signer type.

\begin{table}[t]
\centering
\caption{Cosine retrieval accuracy stratified by signer type
(leave-one-out, $N=64$, Telegram dataset).}
\label{tab:cosine_detail}
\begin{tabular}{@{}lccc@{}}
\toprule
Metric & All ($n{=}335$) & Translator ($n{=}119$) & User ($n{=}216$) \\
\midrule
Top-1 Accuracy (\%) & 30.7 & 25.2 & 33.8 \\
Top-3 Accuracy (\%) & 44.2 & 37.8 & 47.7 \\
\bottomrule
\end{tabular}
\end{table}

User videos achieve higher accuracy than translator videos,
consistent with the Seq2Seq results.
User videos outnumber translator videos approximately 2:1, providing
denser coverage in the embedding space.
Consequently, user queries are more likely to find close neighbours
within the same class.

% ----------------------------------------------------------------
\subsection{Seq2Seq: Training Dynamics}\label{sec:training_results}
% ----------------------------------------------------------------

Table~\ref{tab:training} summarises the Seq2Seq training trajectory
on the Telegram dataset.
The model was trained for 71~epochs with early stopping (patience~7).

\begin{table}[t]
\centering
\caption{Seq2Seq training dynamics on the Telegram dataset (selected epochs).}
\label{tab:training}
\begin{tabular}{@{}cccc@{}}
\toprule
Epoch & Train Loss & Val Loss & Learning Rate \\
\midrule
1  & 5.752 & 4.894 & 0.0100 \\
5  & 3.505 & 3.973 & 0.0100 \\
10 & 3.169 & 3.880 & 0.0100 \\
20 & 1.663 & 3.128 & 0.0100 \\
30 & 1.173 & 3.029 & 0.0050 \\
40 & 1.178 & 2.801 & 0.0050 \\
50 & 0.744 & 2.678 & 0.0025 \\
60 & 0.730 & 2.620 & 0.0025 \\
64 & 0.689 & \textbf{2.568} & 0.0025 \\
71 & 0.626 & 2.595 & 0.00125 \\
\bottomrule
\end{tabular}
\end{table}

The best validation loss (2.568) is achieved at epoch~64, after which
validation loss begins to increase while training loss continues to
decrease, indicating overfitting despite regularisation.
Early stopping restores the epoch-64 checkpoint.
The gap between train and validation loss ($\Delta = 1.97$ by
epoch~71) reflects the limited size of the training corpus: with
only 335~sentence-level videos spread across 120~unique sentences,
the model memorises training patterns rather than learning fully
generalisable sequence mappings.
The learning rate is reduced four times (0.01 $\rightarrow$ 0.005
$\rightarrow$ 0.0025 $\rightarrow$ 0.00125) via
ReduceLROnPlateau, each reduction yielding a brief improvement
in validation loss.

% ----------------------------------------------------------------
\subsection{Effect of Frame Count}\label{sec:framecount}
% ----------------------------------------------------------------

We evaluate the sensitivity of the temporal sampling parameter
$N \in \{16, 32, 64, 128\}$ for the cosine retrieval approach by
resampling stored 64-frame matrices via uniform index selection.
Table~\ref{tab:frames} reports the results.

\begin{table}[t]
\centering
\caption{Cosine retrieval accuracy vs.\ sampled frame count~$N$.}
\label{tab:frames}
\begin{tabular}{@{}ccccc@{}}
\toprule
$N$ & Top-1 (\%) & Top-3 (\%) & Top-5 (\%) & Embedding dim \\
\midrule
16  & 27.8 & 39.5 & 46.8 & 2{,}016 \\
32  & 31.6 & 43.9 & 49.7 & 4{,}032 \\
64  & 30.7 & 45.0 & 49.4 & 8{,}064 \\
128 & 30.7 & 45.0 & 49.4 & 16{,}128 \\
\bottomrule
\end{tabular}
\end{table}

Accuracy improves from $N = 16$ to $N = 32$ as finer temporal
resolution captures hand-shape transitions more effectively.
Performance saturates at $N = 64$; doubling to $N = 128$ yields no
additional gain while doubling the embedding dimensionality.
We therefore adopt $N = 64$ as the operating point, offering the best
trade-off between accuracy and computational cost.

% ----------------------------------------------------------------
\subsection{Latency Analysis}\label{sec:latency}
% ----------------------------------------------------------------

Table~\ref{tab:latency} reports average inference latency per query
for both approaches.

\begin{table}[t]
\centering
\caption{Latency breakdown (average per query).}
\label{tab:latency}
\begin{tabular}{@{}lcc@{}}
\toprule
Stage & Cosine (ms) & Seq2Seq (ms) \\
\midrule
Frame read + sampling (from disk) & 667.8 & --- \\
MediaPipe extraction (64 frames)  & 2{,}036.2 & --- \\
Feature loading (cached)          & ---  & 0.9 \\
Cosine search / Beam decode       & 2.8  & 17.1 \\
\midrule
Matching latency (post-feature)   & \textbf{2.8} & 18.0 \\
Total (from video file)           & 2{,}706.8 & 18.0$^*$ \\
\bottomrule
\end{tabular}
\par\smallskip
\footnotesize{$^*$Seq2Seq total assumes pre-extracted cached features.
Full from-video extraction (SqueezeNet + LSTM encoding) adds
$\sim$2\,s per video.}
\end{table}

For the cosine retrieval pipeline, the dominant cost is MediaPipe
landmark extraction ($\approx$75\% of total time).
In the real-time webcam system, extraction occurs incrementally during
recording at approximately 30\,FPS, eliminating post-capture delay.
After recording, only the cosine search is executed, completing in under
3\,ms for all 335 reference vectors.

The Seq2Seq pipeline's matching latency (beam decode) is 17.1\,ms—still
real-time but 6$\times$ slower than cosine search.
When features are pre-cached, the total Seq2Seq inference completes in
18\,ms.
For both approaches, perceived user latency is dominated by the
2-second cooldown confirmation window rather than computation.

% ----------------------------------------------------------------
\subsection{Discussion}\label{sec:discussion}
% ----------------------------------------------------------------

\paragraph{Seq2Seq performance on crowdsourced data.}
The Seq2Seq model achieves 35.5\% Top-1 accuracy on the Telegram
dataset, outperforming the training-free cosine retrieval baseline
(30.7\%).
While modest in absolute terms, this result is notable given the
extremely small training corpus (335~videos, 120~sentences) and the
high variability inherent in crowdsourced recordings.
The BLEU-4 score of 31.43\% indicates that even incorrect
predictions share substantial lexical overlap with ground-truth
sentences.

\paragraph{Signer-type effects.}
User videos achieve higher accuracy (37.0\% vs.\ 32.8\%) across
both pipelines.
This result is attributed to dataset composition: user videos
outnumber translator videos approximately 2:1 (216 vs.\ 119),
providing denser per-sentence coverage.
The per-sentence accuracy analysis reveals that sentences with only
one video achieve 57.7\% average accuracy, while sentences with
exactly two videos drop to 20.7\%, suggesting that intra-class
variance across multiple signers challenges the model's
generalisation.

\paragraph{Mode collapse.}
Analysis of Seq2Seq predictions reveals significant mode collapse:
the most frequent prediction (``m\textipa{\textschwa}n
u\c{s}aql\i{}qdan \c{s}\textipa{\textschwa}h\textipa{\textschwa}r\textipa{\textschwa}
var'') accounts for 36.7\% of all outputs (123/335).
The model generates only 113 unique predictions out of 335 inputs,
and the top-10 predictions cover 52.8\% of all outputs.
This indicates that the decoder converges to statistically common
patterns rather than fully conditioning on the visual encoder output,
a well-documented challenge in low-resource sequence-to-sequence
models.

\paragraph{Vocabulary analysis.}
The model produces 275 unique words, of which 248 (90.2\%) appear in
the ground-truth vocabulary (356 unique words).
Only 27 out-of-vocabulary words are generated, indicating that the
decoder has learned the target lexicon well but struggles with
correct sentence-level composition.
The average predicted sentence length (4.7 words) closely matches
the ground truth (4.9 words).

\paragraph{Strengths of cosine retrieval.}
The proposed retrieval system requires \emph{no training} and
\emph{no GPU}, making it immediately deployable on commodity hardware.
Vocabulary expansion requires only the addition of reference videos.
The cosine search over 335 vectors completes in under 3\,ms.

\paragraph{Strengths of Seq2Seq.}
Despite mode collapse, the Seq2Seq architecture demonstrates that
learned models can outperform retrieval baselines even with very
small crowdsourced datasets.
With more data, it can generate novel sentences not in the reference
database, handle grammatical variation, and generalise across signer
styles through learned representations.

\paragraph{Limitations.}
The primary limitation is dataset size: 335~videos across
120~sentences provides only 2.8 videos per sentence on average.
The Seq2Seq model exhibits clear overfitting (train--val gap of
1.97) and mode collapse (36.7\% of predictions map to a single
sentence).
The cosine retrieval approach is limited by similar hand
configurations across semantically different sentences.
Furthermore, only manual (hand) features are used; non-manual
markers such as facial expressions, head tilt, and body posture—which
carry grammatical meaning in signed languages—are currently excluded.

\paragraph{Future Directions.}
Expanding the crowdsourced dataset via the Telegram bot with
additional signers and sentences is the most direct path to
improved performance.
Integrating MediaPipe Holistic features (face mesh and pose
landmarks) would capture non-manual signals.
For the Seq2Seq model, pre-training on larger sign-language corpora
(e.g.\ PHOENIX-2014T, CSL) followed by fine-tuning on AzSL could
mitigate mode collapse and improve generalisation.
A hybrid approach combining cosine retrieval for candidate generation
and Seq2Seq for re-ranking may combine the strengths of both
approaches.

% ================================================================
\section{Conclusion}\label{sec:conclusion}
% ================================================================

This paper presented and compared two approaches for real-time
Azerbaijani Sign Language (AzSL) sentence recognition, evaluated
exclusively on a crowdsourced Telegram dataset comprising 335~videos
across 120~unique sentences contributed by both professional
translators (119~videos, 35.5\%) and community users
(216~videos, 64.5\%) via a Telegram bot
platform~\cite{mustafazada2025crowdsourcing}.

The first approach—a training-free cosine retrieval pipeline—extracts
MediaPipe hand landmarks from uniformly sampled video frames,
constructs a compact $64 \times 126$ feature matrix per video, and
retrieves the closest matching sentence from a pre-computed database
using cosine similarity.
The second—a Seq2Seq encoder--decoder model—processes video frames
through a SqueezeNet~1.1 CNN backbone and a bidirectional LSTM
encoder (hidden size 256, output dimension 512), then generates
target sentences through an attention-augmented LSTM decoder with
beam search ($B = 5$).
A shared webcam-based inference module automatically detects signing
onset and offset through a four-state finite state machine
(\textsc{idle} $\rightarrow$ \textsc{recording} $\rightarrow$
\textsc{cooldown} $\rightarrow$ \textsc{matching}),
returning the top-3 candidate sentences in real time on standard
CPU hardware.

\paragraph{Key findings.}
On the Telegram dataset, the Seq2Seq model achieved 35.5\% top-1
accuracy (119/335 correct), outperforming the training-free cosine
retrieval baseline at 30.7\% top-1.
At the sequence level, the Seq2Seq model attained a corpus-average
BLEU-4 of 31.43\%, a Word Error Rate (WER) of 64.05\%, and a
Character Error Rate (CER) of 50.64\%.
Both approaches maintain sub-20\,ms matching latency on a standard
CPU, confirming real-time feasibility.

\paragraph{Signer-type analysis.}
User-contributed videos achieved 37.0\% top-1 accuracy compared to
32.8\% for translator videos, with consistent trends across all
metrics (User BLEU-4: 31.62\% vs.\ Translator: 31.09\%; User WER:
63.00\% vs.\ Translator: 65.94\%).
This demonstrates that crowdsourced data diversity benefits learned
models despite increased signer variability, and that the larger
pool of user recordings (216 vs.\ 119) provides denser per-sentence
coverage for more robust generalisation.

\paragraph{Training dynamics.}
The Seq2Seq model trained for 71~epochs before early stopping, with
the best validation loss (2.568) at epoch~64.
The learning rate was reduced four times via ReduceLROnPlateau
($0.01 \rightarrow 0.005 \rightarrow 0.0025 \rightarrow 0.00125$),
each reduction producing a brief improvement in validation loss.
The widening train--validation gap ($\Delta = 1.97$ at epoch~71)
indicates overfitting, a direct consequence of the limited corpus
size (2.8~videos per sentence on average).

\paragraph{Challenges.}
Mode collapse remains the most significant obstacle: the single most
frequent prediction accounts for 36.7\% of all Seq2Seq outputs
(123/335), and the top-10 predictions cover 52.8\% of all outputs.
The model generates only 113~unique predictions out of 335~inputs,
indicating that the decoder defaults to statistically dominant
patterns rather than fully conditioning on the visual encoder
output.
Nevertheless, vocabulary quality is encouraging: the model produces
275~unique words, of which 248 (90.2\%) appear in the ground-truth
vocabulary, and average predicted sentence length (4.7~words)
closely matches the reference (4.9~words).

\paragraph{Future work.}
Several directions remain to address these limitations.
First, expanding the crowdsourced dataset via the Telegram bot with
additional signers, sentences, and recording environments would
directly mitigate overfitting and mode collapse.
Second, a hybrid architecture that uses cosine retrieval for
candidate generation followed by Seq2Seq re-ranking may combine the
complementary strengths of both methods—retrieval robustness with
generative flexibility.
Third, incorporating body pose and facial expression features from
MediaPipe Holistic would capture non-manual markers (e.g.\ head
tilt, eyebrow raise, mouth shape) that carry grammatical meaning in
sign languages and are currently unmodelled.
Fourth, pre-training the encoder on larger multilingual
sign-language corpora before fine-tuning on AzSL could improve
feature representations and reduce the data requirements.
Finally, porting the system to a mobile application would bring
real-time sign language recognition directly to end users in the
Azerbaijani DHH community, fulfilling the accessibility goal that
motivates this work.

% ================================================================
% Bibliography
% ================================================================
%
% \bibliographystyle{splncs04}
% \bibliography{refs}
%
\begin{thebibliography}{10}

\bibitem{alishzade2024azsld}
Alishzade, N., Hasanov, J.:
AzSLD: Azerbaijani Sign Language Dataset for Fingerspelling, Word, and
Sentence Translation with Baseline Software.
arXiv preprint arXiv:2411.xxxxx (2024)

\bibitem{mustafazada2025crowdsourcing}
Mustafazada, S., Muradova, G., Gozalov, R., Garayev, E., Hasanov, J.:
Crowdsourcing-Based Interactive Learning Platform for the International
Sign Language Community.
In: 2025 6th International Conference on Problems of Cybernetics and
Informatics (PCI), pp.~1--5 (2025).
\doi{10.1109/PCI66488.2025.11219822}

\bibitem{fang2013bayesian}
Fang, X., Dehak, N., Glass, J.:
Bayesian Distance Metric Learning on i-vector for Speaker Verification.
ISCA Archive (2013)

\bibitem{dehak2010cosine}
Dehak, N., Dehak, R., Glass, J., Reynolds, D., Kenny, P.:
Cosine Similarity Scoring without Score Normalization Techniques.
Spoken Language Systems Group, MIT (2010)

\bibitem{zhang2020deep}
Zhang, D., Li, Y., Zhang, Z.:
Deep Metric Learning with Spherical Embedding.
In: NeurIPS, vol.~\textbf{33}, pp.~18772--18783 (2020)

\bibitem{zhu2020orthogonality}
Zhu, Y., Mak, B.:
Orthogonality Regularizations for End-to-End Speaker Verification.
ISCA Archive (2020)

\bibitem{liu2024skeleton}
Liu, L., Zheng, H., Zhu, Z., Zhou, P.:
Skeleton-based Sign Language Recognition Using a Dual-Stream
Spatio-Temporal Dynamic Graph Convolutional Network.
arXiv preprint arXiv:2402.xxxxx (2024)

\bibitem{draganov2024hidden}
Draganov, A., Vadgama, S., Bekkers, E.J.:
The Hidden Pitfalls of the Cosine Similarity Loss.
arXiv preprint arXiv:2403.xxxxx (2024)

\bibitem{weber2025timepoint}
Weber, R.S., Ishay, S.B., Lavrinenko, A., Finder, S.E., Freifeld, O.:
TimePoint: Accelerated Time Series Alignment via Self-Supervised Keypoint
and Descriptor Learning.
In: Proc.\ ICML 2025, PMLR vol.~\textbf{267} (2025)

\bibitem{camgoz2018neural}
Camg\"{o}z, N.C., Hadfield, S., Koller, O., Ney, H., Bowden, R.:
Neural Sign Language Translation.
In: CVPR, pp.~7784--7793 (2018)

\bibitem{iandola2016squeezenet}
Iandola, F.N., Han, S., Moskewicz, M.W., Ashraf, K., Dally, W.J.,
Keutzer, K.:
SqueezeNet: AlexNet-level Accuracy with 50x Fewer Parameters and
$<$0.5\,MB Model Size.
arXiv preprint arXiv:1602.07360 (2016)

\end{thebibliography}
\end{document}
